%====================================================================
% PACKAGES
%====================================================================
% marge
\usepackage[margin=1in]{geometry}
% langue
\usepackage[french]{babel}
\usepackage[T1]{fontenc}
\usepackage{microtype}
% police
\usepackage{cmbright}
\usepackage{frenchmath}
\usepackage{mathpazo}
%\usepackage[frenchmath]{mathastext}
% math utils
\usepackage{mathtools}
\usepackage{yhmath}
% image
\usepackage{graphicx}
% dessin
\usepackage{tikz}
\usepackage{tkz-fct}
\usepackage{tkz-tab}
\usepackage{tkz-euclide}
% plusieurs colonnes
\usepackage{multicol}
% soulignements différents
\usepackage{ulem}
% tableaux
\usepackage{tabularray}
% personnaliser les listes
\usepackage{enumitem}
% unités
\usepackage{siunitx}
\sisetup{
    locale=FR
}
% exercice
% use-files -> xsim utilise mets chaque exercice dans un fichier séparé
% use-aux -> xsim utilise le fichier aux de latex au lien d'en créer un nouveau
\usepackage[use-files,use-aux]{xsim}
\usepackage{tcolorbox}
\usepackage{stmaryrd}
%
\usepackage{subfiles}
\usepackage{fontawesome5}
\usepackage{tasks}
\faStyle{regular}
%====================================================================
% SETTINGS
%====================================================================
% XSIM
%% XSIM TEMPLATE
\DeclareExerciseEnvironmentTemplate{math-exam}%
{%
    \GetExercisePropertyTF{points}%
    % s'il y a des points
    {
        \colorbox{black}{%
            \bfseries\color{white}
            \XSIMmixedcase{\GetExerciseName}~\GetExerciseProperty{counter}}
        \hfill
        \colorbox{black}{%
            \bfseries\color{white}%
            \GetExerciseProperty{points}
            \XSIMtranslate{point-abbr}}
        \par
    }
    % sinon
    {
        \colorbox{black}{%
                    \bfseries\color{white}%
                    \XSIMmixedcase{\GetExerciseName}~\GetExerciseProperty{counter}%
                }
        \par
    }%
}%
{%
    \par%
} % TODO: à changer par la commande de xsim appropriée
\xsimsetup{
    path=xsim,% stocker les fichiers auxiliaires dans ce dossier
    %solution/print=true, % afficher les solutions en dessous des questions
    blank/style=\dotuline{#1},
    exercise/template=math-exam,
    solution/template=math-exam
}

% ENUMITEM
% mettre les labels sous forme de lettre en gras
\setlist[enumerate]{
    label=\textbf{\alph*)},
    itemsep=1ex}
% nouveau paramètre pour mettre les listes sur plusieurs colonnes
\SetEnumitemKey{cols}{%
    itemsep = 1\itemsep,
    parsep = 1\parsep,
    before = \raggedcolumns\begin{multicols}{#1},
        after = \end{multicols}}
% \newlist{qcm}{itemize}{1}
% \setlist[qcm]{label=\faSquare, cols=2}
% TIKZ
\usetikzlibrary{math,babel,patterns,calc}
% TASKS
\NewTasksEnvironment[label=\faSquare, label-width=12pt]{QCM}[\choice](2)
\settasks{
    label=\textbf{\alph*)},
    label-width=12pt
}
%============================================
% PERSONNAL MACRO
%============================================
\newtcolorbox{consigne}{colback=gray!5!white,
colframe=gray!75!black}